\section{Conservation as an independent cross-check}
\label{sec:conservation-cross-check}
An improvement in a chosen reflection residual need not improve the magnetic
model as a whole. Conservation supplies a separate test: a scalar co-energy
requires reciprocal differential cross-inductances and zero circulation in
the appropriate current coordinates \cite{haus1989,jebai2014}.
The normalized field $\psi_d=1+y^2$, $\psi_q=y$ satisfies the expected
parity but has curl $-2y$. Conversely, the exact coherent angle rotation in
\cref{eq:combined-exact} preserves integrability and rotates the symmetry
axis. These constraints therefore test different aspects of the baseline.

For the reconstructed field define
$\fluxcurlresidual=\partial_{i_d}\psi_q-\partial_{i_q}\psi_d$.
The dq integrals $\pathintegralA$ and $\pathintegralB$ follow the d-then-q
and q-then-d paths from the origin to $(x,y)=(i_d,i_q)$. With oriented limits,
\begin{align}
 \pathintegralA-\pathintegralB
 &=\int_0^x\int_0^y r_{\rm curl}(\xi,\eta)\,d\eta\,d\xi,\\
 \pathresistanceerror
 &=\frac{\omegae(W_A-W_B)}{2xy}
 =\frac{\omegae}{2|xy|}\iint_{\mathcal A_{\rm geom}}r_{\rm curl}\,dA,
 \label{eq:diagnostic-path-equivalent}\\
 \pathresistance&=\hat R_s-\Delta R_{\rm eq,path}.
\end{align}
The boundary runs counterclockwise for $xy>0$ and clockwise for $xy<0$.
The target indexes the integration rectangle. This is an area-averaged curl
equivalent, not a local resistance at that point. Overlapping rectangles
share samples, and signed curl can cancel over their areas.
It equals the resistance-mismatch contribution only under the sole-mismatch
baseline. A voltage error proportional to current remains exactly confounded.

The existing audit uses measured-current Delaunay interpolation of each flux
component and excludes targets below ten percent of either axis maximum or
with unsupported paths. \Cref{tab:diagnostic-path-comparison} compares the
result with the high-current symmetry summary. Accepted rectangle counts are
2868/2869/2839 in the thermal slices and 65/45 in the speed slices. Hull skips
after the axis threshold are 3/2/32 and 4/24, respectively. These sets differ
from the pair support in \cref{tab:experimental-summary}; the two speed path
medians are not a pointwise speed comparison. The separate direct test on
36 common valid pairs remains the relevant matched-speed check.
\begin{table}[tb]
\centering\small
% Generated by numerics/export_audit_tables.py; do not edit.
% audit_results.json SHA256 1523aedbba70a4364963d65cac529a522b4208336296fdc04ce6b2cdff27a765
% observable_checks.json SHA256 f1235618a0bf5fe74cc0a84385a49cb39d5da9a8db0b322ebd75022387a00d54
\begin{tabular}{rrrrrrr}
\toprule
$n$ & $T_{\rm rot,ref}$ & $N_{\rm rect}$ & $R_{\rm eq,sym}$ & $R_{\rm eq,path}$ & MAD & P1 exact \\
\midrule
2000 & 40 & 2868 & 7.599140 & 5.007865 & 1.783325 & 5.007792 \\
2000 & 60 & 2869 & 8.274010 & 5.885553 & 1.728698 & 5.885512 \\
2000 & 80 & 2839 & 8.963406 & 6.488047 & 1.769623 & 6.487975 \\
1000 & 70 & 65 & 12.156625 & 5.960078 & 2.405193 & 5.960063 \\
3000 & 70 & 45 & 15.334054 & 3.809622 & 2.991105 & 3.808417 \\
\bottomrule
\end{tabular}

\caption{Conservation cross-check on unchanged audit support. All resistance
columns are m$\Omega$; $n$ is rpm and rotor reference is $^\circ$C.
Symmetry uses the high-current subset; the path columns are median/raw MAD
over $N_{\rm rect}$ rectangles. P1 exact means exact integration of the same
interpolator, not exact physical identification.}
\label{tab:diagnostic-path-comparison}
\end{table}

At 2000 rpm/60 $^\circ$C, the conditional symmetry value $\symmetryresistance$
\SI{8.274010}{\milli\ohm} differs from the path median
\SI{5.885553}{\milli\ohm}. Exact P1 integration changes the latter only to
\SI{5.885512}{\milli\ohm}. The 500-point trapez error over all 2869 targets
is \SI{0.0000766}{\milli\ohm} RMS, far below the discrepancy. Exact path
integration and independent triangle-area Green evaluation agree within
$2.54\cdot10^{-14}$ Wb A on eight checked rectangles.

Changing the used resistance from \SI{7.695}{\milli\ohm} to the symmetry
equivalent reduces selected high-current parity, but increases all-pair d
and vector residual RMS and the median relative path difference
\begin{equation}
 \relativepathdifference=|W_A-W_B|/[(|W_A|+|W_B|)/2].
\end{equation}
This ratio is used where its denominator is nonzero; it is not relative error
against a measured magnetic potential. \Cref{tab:diagnostic-adjustment} retains
the support and both residual
components. No measured parameter correction is thereby validated.
\begin{table}[tb]
\centering\small
% Generated by numerics/export_audit_tables.py; do not edit.
% audit_results.json SHA256 1523aedbba70a4364963d65cac529a522b4208336296fdc04ce6b2cdff27a765
% observable_checks.json SHA256 f1235618a0bf5fe74cc0a84385a49cb39d5da9a8db0b322ebd75022387a00d54
\begin{tabular}{lllrr}
\toprule
Support & Metric & Unit & Raw & Adjusted \\
\midrule
High (80 pairs) & $r_d$ RMS & mWb & 1.100092 & 0.219939 \\
High (80 pairs) & $r_q$ RMS & mWb & 0.487494 & 0.277063 \\
High (80 pairs) & Vector RMS & mWb & 1.203268 & 0.353747 \\
All (1172 pairs) & $r_d$ RMS & mWb & 1.746480 & 1.954456 \\
All (1172 pairs) & $r_q$ RMS & mWb & 1.308799 & 1.183835 \\
All (1172 pairs) & Vector RMS & mWb & 2.182464 & 2.285030 \\
2869 rectangles & Median $\rho_{\rm path}$ & \% & 3.488090 & 4.226703 \\
2869 rectangles & Median $\Delta R_{\rm eq,path}$ & m$\Omega$ & 1.809447 & 2.388457 \\
2869 rectangles & Median $R_{\rm eq,path}$ & m$\Omega$ & 5.885553 & 5.885553 \\
\bottomrule
\end{tabular}

\caption{Raw versus symmetry-adjusted resistance hypothesis at
2000 rpm/60 $^\circ$C, on unchanged data. High-current improvement accompanies
worse global d/vector parity and relative path disagreement.}
\label{tab:diagnostic-adjustment}
\end{table}
The pointwise invariance of $R_{\rm eq,path}$ under any constant used-resistance
shift is algebraic: reconstruction adds $h\matr J\vect i/\omega_e$, whose
path equivalent is $h$. Thus $(R+h)-\Delta R_{\rm eq,path}(R+h)
=R-\Delta R_{\rm eq,path}(R)$. This checks affine interpolation and integration,
and does not calibrate winding resistance.

Numerical and physical limits remain separate. The componentwise P1 field is
$C^0$ with trianglewise gradients that jump at edges; a common mesh enforces
no symmetric Jacobian. The 60-$^\circ$C mesh has very skinny triangles and
extreme local curl equivalents, which are derivative-conditioning diagnostics,
not physical resistances. Sampling an analytic conservative nonlinear field
on the same geometry also creates nonzero interpolated circulation. Exact
quadrature removes integration error without validating the interpolated field.

Nor do terminal-current coordinates necessarily describe a conservative
magnetic storage law when magnetizing and loss currents differ
\cite{richter2014}. The recorded current RMS ratios are consistent with
amplitude-invariant dq, but the transform and voltage/power normalization
are not fully verified. $W_A,W_B$ are reconstructed dq flux integrals in Wb A,
not independently measured absolute magnetic energy; a consistent global
$3/2$ cannot remove their disagreement. The independent conservation test
strengthens the evidence boundary: smaller selected symmetry residuals alone
do not establish a valid automatic resistance correction or identify iron loss.
